%% kapa: exam=2 round=24 year=2013 period=1 subject=appraisal_practice subject_ko=감정평가실무 seq=2 number=2 points=30
%% source: Q-Net 감정평가사 기출문제 (cst00309) artlSeq=5077404 file=extracted/1008/5077404/2013년 제24회 감정평가사 2차시험 기출문제/1교시.hwp
%% conversion: kordoc 4.24.51 -> tools/gen2.py (auto); review=proofread-sonnet; flags=none
\begin{longquestion}{【문제 2】}
공기업인 (주)H전력은 서울특별시 동대문구 Y동 45번지에 소재하는 토지, 건물(이하 “대상물건”이라 함)에 대하여 ① 매각하는 방안, ② 보상을 받는 방안, ③ 분양을 받는 방안을 감정평가사 甲에게 검토 요청하였다. \par
대상물건은 재개발 정비구역 내에 소재하고 있으며 관리처분까지 확정되어, 현재 공실상태이고, 종후자산에 대한 분양계약을 체결하면 바로 철거가 진행될 예정이며, 종후자산은 5년 후인 2018. 4. 30자로 입주예정이다. \par
재개발조합에서는 미계약자에게 2013. 7. 31 기준시점으로「공익사업을 위한 토지 등의 취득 및 보상에 관한 법률」에 의거 평가한 후 2013. 8. 31자로 일시불로 지급하기로 되어있다. 또한, 대상물건은 2013. 4. 30에 일시불로 80억원에 매입하겠다는 매수희망자가 있다. 다음 물음에 답하시오. (30점) \par
\medskip\noindent 물음 1) 2013. 7. 31 기준시점으로 보상평가액 (15점) \par
\medskip\noindent 물음 2) 2018. 4. 30(입주시) 기준시점의 현금정산액을 포함한 종후자산 가치 (5점) \par
\medskip\noindent 물음 3) (주) H전력이 제시한 3가지 방안을 2013. 4. 30 기준으로 비교 검토한 후 적절한 방안을 제시하고 그 이유를 설명하시오. (10점) \par
\textbf{<자료 1> 대상물건 개요} \par
1. 대상토지 \par
1) 소재지: 서울특별시 동대문구 Y동 45 번지 \par
2) 지목 및 면적: 대, 820m\textsuperscript{2} \par
3) 용도지역: 일반상업지역, 도시계획 도로 일부저촉(20\%) \par
4) 토지의 특성: 장방형, 완경사, 광대소각 \par
2. 대상건물 \par
1) 구조: 철근콘크리트조 평옥개지붕 \par
2) 건물면적: 520m\textsuperscript{2} \par
3) 층수: 지하1층/지상3층 \par
4) 사용승인일: 1980. 9. 7 \par
5) 용도: 업무용 \par
6) 대상건물은 도시계획도로에 저촉되지 않음 \par
\textbf{<자료 2> 현장조사자료} \par
1. 건물 \par
1) 공부면적: 520m\textsuperscript{2} \par
2) 실측면적: 650m\textsuperscript{2} \par
2. 제시외 건물 \par
1) 수위실: 벽돌조 슬래브지붕 단층 10m\textsuperscript{2} \par
수위실은 전체면적이 도시계획도로에 저촉됨 \par
2) 창고: 목조 와즙 단층 60m\textsuperscript{2} \par
3) 제시외 건물들은 1985. 10. 10 신축된 것으로 조사되었음 \par
3. 기타 지장물 \par
1) 벽돌조 담장: 110m\textsuperscript{2} \par
2) 바닥 포장: 아스콘 포장 530m\textsuperscript{2} \par
3) 축대: 철근콘크리트 54m\textsuperscript{2} \par
4) 수목 \par
\medskip\noindent (1) 소나무 45년생 3주 \par
\medskip\noindent (2) 감나무 25년생 5주 \par
\medskip\noindent (3) 대추나무 15년생 5주 \par
\textbf{<자료 3> 재개발관련 자료} \par
1. 사업승인일: 2012. 4. 5 \par
2. 종전자산 평가액 \par
1) 토지: 7,790,000,000원 (820m\textsuperscript{2} × 9,500,000원/m\textsuperscript{2}) \par
2) 건물: 166,400,000원 (520m\textsuperscript{2} × 320,000원/m\textsuperscript{2}) \par
3) 합계: 7,956,400,000원 \par
4) 비례율: 95\% \par
3. 종후자산 \par
1) (주)H전력은 상가 1층 대지권 290m\textsuperscript{2}, 건물(전용면적) 720m\textsuperscript{2}을 분양가능 \par
2) 종전자산보다 종후자산이 많을 경우에는 입주시 일시불로 차액을 지급하고, 적을 경우에는 입주시(2018. 4. 30) 일시불로 차액을 받을 수 있음 \par
3) 상가의 조합원 분양가격은 건물면적(전용면적)기준 1층 @8,000,000원/m\textsuperscript{2}, 2층 @3,000,000원/m\textsuperscript{2}, 3층 @1,800,000원/m\textsuperscript{2}, 지층 @1,300,000원/m\textsuperscript{2} 임 \par
\textbf{<자료 4> 공시지가 표준지와 지가변동률} \par
1. 인근 표준지공시지가 \par
{\small\centering
\begin{tabular}{|c|c|c|c|c|c|c|c|c|c|}
\hline
\multirow{2}{*}{\makecell{기호}} & \multirow{2}{*}{\makecell{소재지}} & \multirow{2}{*}{\makecell{면적\newline{}(m\textsuperscript{2})}} & \multirow{2}{*}{\makecell{지목}} & \multirow{2}{*}{\makecell{이용\newline{}상황}} & \multirow{2}{*}{\makecell{용도지역}} & \multirow{2}{*}{\makecell{도로교통}} & \multirow{2}{*}{\makecell{형상\newline{}지세}} & \multicolumn{2}{|c|}{\makecell{공시지가(원/m\textsuperscript{2})}} \\
\cline{9-10}
 &  &  &  &  &  &  &  & \makecell{2012년} & \makecell{2013년} \\
\hline
\makecell{①} & \makecell{Y동\newline{}26-5} & \makecell{810} & \makecell{대} & \makecell{상업용} & \makecell{일반상업} & \makecell{광대소각} & \makecell{세장형 평지} & \makecell{6,000,000} & \makecell{6,500,000} \\
\hline
\makecell{②} & \makecell{Y동 32-1} & \makecell{302} & \makecell{대} & \makecell{업무용} & \makecell{일반상업} & \makecell{세로\newline{}(가)} & \makecell{세장형\newline{}평지} & \makecell{4,700,000} & \makecell{4,900,000} \\
\hline
\end{tabular}\par}

※ 기호② 표준지는 도시계획도로에 30\% 저촉됨 \par
2. 시점수정 관련자료 \par
1) 지가변동률(동대문구 상업지역) \par
{\small\centering
\begin{tabular}{|c|c|c|c|c|}
\hline
\makecell{2012.1.1～3.31} & \makecell{2012년 4월} & \makecell{2012년 누계} & \makecell{2013.1.1～3.31} & \makecell{2013년 3월} \\
\hline
\makecell{0.572\%} & \makecell{0.180\%} & \makecell{2.240\%} & \makecell{0.512\%} & \makecell{0.155\%} \\
\hline
\end{tabular}\par}

※ 2013년 4월 이후 지가변동률은 미고시 되었으므로 2013년 3월 지가변동률을 연장 적용하기로 함 \par
2) 동대문구 Y동 일대의 상가는 부동산 경기침체로 2013년 상반기 중 가격변동이 거의 없으므로 시점수정은 없음 \par
3) 2013. 4. 30～2018. 4. 30 사이 5년간의 상가 상승률은 연 2\%로 가정함 \par
\textbf{<자료 5> 가격자료} \par
1. 평가선례 및 거래사례 \par
1) 평가선례 \par
(단위: 원/m\textsuperscript{2}) \par
{\small\centering
\begin{tabular}{|c|c|c|c|c|c|c|c|c|}
\hline
\makecell{기호} & \makecell{소재지} & \makecell{지목} & \makecell{면적\newline{}(m\textsuperscript{2})} & \makecell{이용상황} & \makecell{용도\newline{}지역} & \makecell{도로\newline{}접면} & \makecell{기준시점\newline{}평가목적} & \makecell{평균단가} \\
\hline
\makecell{A} & \makecell{Y동 24-2} & \makecell{대} & \makecell{330} & \makecell{상업용} & \makecell{일반\newline{}상업} & \makecell{광대\newline{}소각} & \makecell{2013. 3. 5\newline{}보상} & \makecell{10,500,000} \\
\hline
\makecell{B} & \makecell{Y동 30-1} & \makecell{대} & \makecell{255} & \makecell{주거용} & \makecell{일반\newline{}상업} & \makecell{세로\newline{}(불)} & \makecell{2012. 2. 10\newline{}담보} & \makecell{7,500,000} \\
\hline
\makecell{C} & \makecell{Y동 24-5} & \makecell{대} & \makecell{420} & \makecell{상업용} & \makecell{일반\newline{}상업} & \makecell{광대\newline{}소각} & \makecell{2012. 4. 2\newline{}보상} & \makecell{9,500,000} \\
\hline
\end{tabular}\par}

※ 기호 A와 C는 면적의 40\%가 도시계획도로에 저촉됨 \par
2) 거래사례 \par
(단위: 원/m\textsuperscript{2}) \par
{\small\centering
\setlength{\tabcolsep}{3pt}
\begin{tabular}{|c|c|c|c|c|c|c|c|c|c|}
\hline
\makecell{기호} & \makecell{소재지} & \makecell{지목} & \makecell{면적\newline{}(m\textsuperscript{2})} & \makecell{이용\newline{}상황} & \makecell{용도\newline{}지역} & \makecell{도로\newline{}접면} & \makecell{거래일자} & \parbox{1.8cm}{\centering 토지기준거래단가} & \makecell{비고} \\
\hline
\makecell{D} & \makecell{Y동 42-7} & \makecell{대} & \makecell{550} & \makecell{상업용} & \makecell{일반\newline{}상업} & \makecell{광대\newline{}소각} & \makecell{2013. 3. 27} & \makecell{12,000,000} & \parbox{1.8cm}{\centering 토지단가에 건물가포함} \\
\hline
\makecell{E} & \makecell{Y동 30-7} & \makecell{대} & \makecell{187} & \makecell{주상용} & \makecell{일반\newline{}상업} & \makecell{세로\newline{}(가)} & \makecell{2013. 2. 15} & \makecell{8,000,000} & \makecell{토지단가에\newline{}건물가포함} \\
\hline
\end{tabular}\par}

※ 기호 E토지는 지상건물을 포함하여 거래하였으나 노후 건물로 철거예정이며 철거비와 폐자재판매비가 거의 유사함 \par
3) 상가거래사례 \par
(단위: 천원) \par
{\small\centering
\setlength{\tabcolsep}{3pt}
\begin{tabular}{|c|c|c|c|c|c|c|c|c|}
\hline
\makecell{기호} & \makecell{소재지} & \makecell{면적\newline{}(m\textsuperscript{2})} & \makecell{이용\newline{}상황} & \makecell{용도\newline{}지역} & \makecell{도로\newline{}접면} & \makecell{거래일자} & \makecell{거래금액} & \makecell{비고} \\
\hline
\makecell{F} & \makecell{Y동 22} & \makecell{건물(전용) 640\newline{}대지권 160} & \makecell{상업용} & \makecell{일반\newline{}상업} & \makecell{광대\newline{}소각} & \makecell{2013. 1. 30} & \makecell{7,040,000} & \makecell{OO단지내\newline{}상가 1층} \\
\hline
\makecell{G} & \makecell{Y동 22} & \makecell{건물(전용) 530\newline{}대지권 132} & \makecell{상업용} & \makecell{일반\newline{}상업} & \makecell{광대\newline{}소각} & \makecell{2013. 2. 15} & \makecell{1,060,000} & \makecell{OO단지내 상가 3층} \\
\hline
\end{tabular}\par}

※ OO단지는 2012년 신축하여 입주한 아파트단지 내 상가로 대상물건이 소재한 정비구역의 재개발이 완료되면 제반 단지여건이 유사함 \par
4) 기타 가격자료 \par
\medskip\noindent (1) 동 지역은 재개발사업의 사업승인 후 지가가 상당한 폭으로 상승하였음 \par
\medskip\noindent (2) 상가는 조합원 분양가보다 일반분양가가 약 45\% 높게 결정되었으며 70\% 정도 분양되었음 \par
2. 대상건물 및 지장물 가격자료 \par
1) 대상건물 \par
\medskip\noindent (1) 재조달원가: 보정단가 포함한 전체면적 기준 @950,000원/m\textsuperscript{2} \par
\medskip\noindent (2) 내용년수: 철근콘크리트조이므로 50년으로 함 \par
2) 제시외 건물 \par
\medskip\noindent (1) 재조달원가: 수위실 @320,000원/m\textsuperscript{2}, 창고 @600,000원/m\textsuperscript{2} \par
\medskip\noindent (2) 내용년수: 수위실, 창고 모두 40년으로 함 \par
3) 지장물 가격자료 \par
감가상각 후 적용단가임 \par
\medskip\noindent (1) 벽돌조 담장: @40,000/m\textsuperscript{2} \par
\medskip\noindent (2) 아스콘포장 비용: @90,000원/m\textsuperscript{2} \par
\medskip\noindent (3) 축대 조성비용: @120,000원/m\textsuperscript{2} \par
\medskip\noindent (4) 수목가격자료 \par
{\small\centering
\begin{tabular}{|c|c|c|c|}
\hline
\makecell{수종} & \makecell{단위} & \makecell{취득비} & \makecell{이식비} \\
\hline
\makecell{소나무 45년생} & \makecell{1주} & \makecell{15,000,000원} & \makecell{4,200,000원} \\
\hline
\makecell{감나무 25년생} & \makecell{1주} & \makecell{500,000원} & \makecell{800,000원} \\
\hline
\makecell{대추나무 15년생} & \makecell{1주} & \makecell{200,000원} & \makecell{180,000원} \\
\hline
\end{tabular}\par}

※ 소나무는 고손율이 20\%이고, 감나무와 대추나무는 고손율이 15\%임 \par
\textbf{<자료 6> 지역요인, 개별요인 비교자료} \par
1. 지역요인 \par
공시지가 비교표준지, 대상토지, 평가선례 및 거래사례는 동일지역임 \par
2. 개별요인비교 \par
1) 토지 개별요인비교 \par
{\small\centering
\begin{tabular}{|c|c|c|c|c|c|c|c|}
\hline
\makecell{대상\newline{}토지} & \makecell{표준지\newline{}①} & \makecell{표준지\newline{}②} & \makecell{평가\newline{}선례 A} & \makecell{평가\newline{}선례 B} & \makecell{평가\newline{}선례 C} & \makecell{거래\newline{}사례D} & \makecell{거래\newline{}사례 E} \\
\hline
\makecell{100} & \makecell{100} & \makecell{75} & \makecell{103} & \makecell{74} & \makecell{102} & \makecell{105} & \makecell{79} \\
\hline
\end{tabular}\par}

2) 상가 개별요인비교 \par
{\small\centering
\begin{tabular}{|c|c|c|}
\hline
\makecell{대상 종후상가} & \makecell{거래사례 F} & \makecell{거래사례 G} \\
\hline
\makecell{100} & \makecell{95} & \makecell{20} \\
\hline
\end{tabular}\par}

3) 공법상 제한사항 \par
{\small\centering
\begin{tabular}{|c|c|c|}
\hline
\makecell{구분} & \makecell{일반} & \makecell{제한} \\
\hline
\makecell{일반} & \makecell{1.0} & \makecell{0.85} \\
\hline
\makecell{제한} & \makecell{1.18} & \makecell{1.0} \\
\hline
\end{tabular}\par}

\textbf{<자료 7> 기타자료} \par
1. 일시불 현가화시 할인율은 연 6\%(월 0.5\%)를 적용함 \par
\end{longquestion}
